\documentclass[12pt,a4paper]{article}
\usepackage[margin=2.5cm]{geometry}
\usepackage{graphicx}
\usepackage{hyperref}
\usepackage{listings}
\usepackage{xcolor}
\usepackage{booktabs}
\usepackage{caption}
\usepackage{url}

\definecolor{codegreen}{rgb}{0,0.6,0}
\definecolor{codegray}{rgb}{0.5,0.5,0.5}
\definecolor{codepurple}{rgb}{0.58,0,0.82}
\definecolor{backcolour}{rgb}{0.95,0.95,0.92}

\lstdefinestyle{mystyle}{
    backgroundcolor=\color{backcolour},
    commentstyle=\color{codegreen},
    keywordstyle=\color{magenta},
    numberstyle=\tiny\color{codegray},
    stringstyle=\color{codepurple},
    basicstyle=\ttfamily\small,
    breaklines=true,
    frame=single,
    numbers=left,
    tabsize=4
}

\lstset{style=mystyle}

\hypersetup{
    colorlinks=true,
    linkcolor=blue,
    filecolor=magenta,
    urlcolor=blue,
}

\title{Traffic Management System: Real-Time Event Processing with PostgreSQL}
\author{Dheeraj Murthy}
\date{\today}

\begin{document}
\maketitle

\begin{abstract}
This report documents a real-time traffic management system built using PostgreSQL LISTEN/NOTIFY for event-driven communication. The system consists of three components: a C++ event watcher that listens for database notifications, a Go REST API server for the traffic dashboard, and PostgreSQL triggers for reactive traffic logic. The system parses ECA (Event-Condition-Action) rules from XML and processes sensor and signal events in real-time.
\end{abstract}

\tableofcontents
\newpage

\section{Introduction}

Modern traffic management systems require real-time processing of sensor data and signal updates. This project implements an event-driven architecture using PostgreSQL's LISTEN/NOTIFY mechanism for immediate notification when database state changes occur.

\subsection{Objectives}
\begin{itemize}
    \item Real-time sensor data processing (vehicle count, average speed)
    \item Signal phase management and timing control
    \item Junction event handling (phase changes, queue overflow, emergencies)
    \item Corridor-based traffic flow optimization
    \item ECA rule engine for automated responses
\end{itemize}

\section{System Architecture}

\subsection{Component Overview}

\begin{figure}[H]
\centering
\fbox{
\parbox{0.8\textwidth}{
\textbf{Architecture Flow:}
\newline\newline
PostgreSQL Database $\rightarrow$ C++ Watcher $\rightarrow$ Event Handlers
\newline\newline
C++ Watcher $\leftrightarrow$ Go API Server
}}
\end{figure}
\bigskip

\begin{table}[H]
\centering
\begin{tabular}{lll} \toprule
Component & Language & Purpose \\ \midrule
Watcher & C++20 & PostgreSQL listener, event dispatcher \\
Server & Go 1.25 & REST API for dashboard \\
Database & PostgreSQL & Schema, triggers, reactive logic \\ \bottomrule
\end{tabular}
\caption{Component Summary}
\end{table}

\subsection{Directory Structure}

\begin{lstlisting}[language=bash,style=mystyle]
project/
├── AGENTS.md           # Project knowledge base
├── server/            # Go API server
│   └── main.go
├── sql/              # PostgreSQL schema
│   ├── db_create.sql
│   ├── db_data.sql
│   └── xml/
│       └── eca-rules-example.xml
└── watcher/          # C++ event watcher
    ├── CMakeLists.txt
    ├── include/watcher/
    │   ├── config.hpp
    │   ├── connection.hpp
    │   ├── dispatcher.hpp
    │   ├── listener.hpp
    │   └── rules/
    │       └── rule_loader.hpp
    └── src/
        ├── main.cpp
        ├── config.cpp
        ├── connection.cpp
        ├── listener.cpp
        ├── dispatcher.cpp
        ├── rules/
        │   └── rule_loader.cpp
        └── handlers/
            ├── sensor_handler.cpp
            └── signal_handler.cpp
\end{lstlisting}

\section{Database Layer}

PostgreSQL serves as the event source with triggers emitting NOTIFY events.

\subsection{Event Channels}

\begin{table}[H]
\centering
\begin{tabular}{ll} \toprule
Channel & Event Type \\ \midrule
sensor\_update\_channel & Vehicle count, speed data \\
signal\_update\_channel & Phase changes, timing \\
junction\_event\_channel & Junction-specific events \\ \bottomrule
\end{tabular}
\end{table}

Sample trigger:
\begin{lstlisting}[language=SQL,style=mystyle]
CREATE TRIGGER sensor_update_trigger
AFTER INSERT ON sensor_data
FOR EACH ROW
EXECUTE FUNCTION notify_sensor_update();
\end{lstlisting}

\section{Event Watcher (C++)}

The C++ watcher connects to PostgreSQL and listens for NOTIFY events.

\subsection{Build System (CMake)}

\begin{lstlisting}[language=CMake,style=mystyle]
cmake_minimum_required(VERSION 3.16)
project(sql_watcher VERSION 1.0.0 LANGUAGES CXX)

set(CMAKE_CXX_STANDARD 20)
set(CMAKE_CXX_STANDARD_REQUIRED ON)

set(PostgreSQL_ROOT "/opt/homebrew/opt/postgresql@17")
find_package(PostgreSQL REQUIRED)

include(FetchContent)
FetchContent_Declare(json GIT_REPOSITORY https://github.com/nlohmann/json.git ...)
FetchContent_MakeAvailable(json)

FetchContent_Declare(pugixml ...)
FetchContent_MakeAvailable(pugixml)

add_executable(watcher ${SOURCES})
target_link_libraries(watcher PostgreSQL::PostgreSQL nlohmann_json pugixml pthread)
\end{lstlisting}

\subsection{Main Entry Point}

\begin{lstlisting}[language=C++,style=mystyle]
#include "watcher/config.hpp"
#include "watcher/connection.hpp"
#include "watcher/dispatcher.hpp"
#include "watcher/handlers/sensor_handler.hpp"
#include "watcher/handlers/signal_handler.hpp"
#include "watcher/listener.hpp"
#include "watcher/rules/rule_loader.hpp"

int main() {
    auto config = watcher::Config::fromEnv();
    config.dbname = "traffic_db";
    config.user = "dheerajmurthy";
    config.password = "dheeraj";

    watcher::Connection conn(config.connectionString());
    std::cout << "Connected.\n";

    watcher::RuleLoader rule_loader;
    rule_loader.load_from_file("eca-rules.xml");
    std::cout << "Loaded " << rule_loader.rules().size() << " rules\n";

    watcher::EventDispatcher dispatcher;
    dispatcher.registerHandler("sensor_update_channel", 
        std::make_unique<watcher::SensorHandler>());
    dispatcher.registerHandler("signal_update_channel", 
        std::make_unique<watcher::SignalHandler>());

    conn.listen("sensor_update_channel");
    conn.listen("signal_update_channel");
    conn.listen("junction_event_channel");

    watcher::Listener listener(conn, dispatcher);
    listener.start();

    std::cout << "Listening for events... (Ctrl+C to exit)\n";
    // Main loop...
}
\end{lstlisting}

\subsection{Event Dispatcher Pattern}

\begin{lstlisting}[language=C++,style=mystyle]
class EventHandler {
public:
    virtual ~EventHandler() = default;
    virtual void handle(const std::string& channel, const json& data) = 0;
    virtual std::string name() const = 0;
};

class EventDispatcher {
public:
    void registerHandler(const std::string& channel, HandlerPtr handler);
    void dispatch(const std::string& channel, const std::string& payload);
};
\end{lstlisting}

Handler implementation:
\begin{lstlisting}[language=C++,style=mystyle]
class SensorHandler : public EventHandler {
public:
    std::string name() const override { return "SensorHandler"; }
    
    void handle(const std::string& channel, const json& data) override {
        std::cout << "Processing sensor update...\n";
        std::cout << "  Data: " << data.dump(2) << "\n";
    }
};
\end{lstlisting}

\section{ECA Rule Engine}

The system loads Event-Condition-Action rules from XML using pugixml.

\subsection{Rule Schema (XSD)}

\begin{lstlisting}[language=XML,style=mystyle]
<?xml version="1.0" encoding="UTF-8"?>
<xsd:schema xmlns:xsd="http://www.w3.org/2001/XMLSchema">
  <xsd:element name="eca-rules" type="tns:EcaRulesType"/>
  <xsd:complexType name="EcaRulesType">
    <xsd:sequence>
      <xsd:element name="rule" type="tns:RuleType" maxOccurs="unbounded"/>
    </xsd:sequence>
    <xsd:attribute name="version" type="xsd:string"/>
    <xsd:attribute name="enabled" type="xsd:boolean"/>
  </xsd:complexType>
  <xsd:complexType name="RuleType">
    <xsd:sequence>
      <xsd:element name="name" type="xsd:string"/>
      <xsd:element name="event" type="tns:EventType"/>
      <xsd:element name="condition" type="tns:ConditionType"/>
      <xsd:element name="action" type="tns:ActionType"/>
    </xsd:sequence>
    <xsd:attribute name="id" type="xsd:ID"/>
    <xsd:attribute name="priority" type="xsd:int"/>
    <xsd:attribute name="enabled" type="xsd:boolean"/>
  </xsd:complexType>
</xsd:schema>
\end{lstlisting}

\subsection{Example Rules}

\begin{lstlisting}[language=XML,style=mystyle]
<?xml version="1.0" encoding="UTF-8"?>
<eca-rules version="1.0" enabled="true">
  <rule id="rule-001" priority="10" enabled="true">
    <name>High Traffic Response</name>
    <description>When queue exceeds 40, extend green time</description>
    <event>
      <sensor-update vehicle-count-min="41"/>
    </event>
    <condition>
      <queue-threshold threshold="40" comparison="gt"/>
    </condition>
    <action>
      <set-phase-timing green-time="90" is-active="true"/>
    </action>
  </rule>
  
  <rule id="rule-002" priority="5" enabled="true">
    <name>Low Traffic Optimization</name>
    <event>
      <sensor-update vehicle-count-max="9"/>
    </event>
    <condition>
      <queue-threshold threshold="10" comparison="lt"/>
    </condition>
    <action>
      <adjust-green-time adjustment-seconds="-10"/>
    </action>
  </rule>
</eca-rules>
\end{lstlisting}

Rule loader header:
\begin{lstlisting}[language=C++,style=mystyle]
enum class EventType {
    SENSOR_UPDATE,
    JUNCTION_EVENT,
    TIME_SCHEDULE,
    CORRIDOR_EVENT,
    UNKNOWN
};

enum class SignalMode {
    DYNAMIC, PERIODIC, MANUAL, ECA_CONTROLLED
};

struct EcaRule {
    std::string id;
    int priority = 0;
    bool enabled = false;
    std::string name;
    EventType event_type;
    std::vector<Action> actions;
};

class RuleLoader {
public:
    bool load_from_file(const std::string& path);
    const std::vector<EcaRule>& rules() const;
    std::vector<const EcaRule*> find_matching_rules(EventType type) const;
};
\end{lstlisting}

\section{API Server (Go)}

\begin{lstlisting}[language=Go,style=mystyle]
package main

import (
    "database/sql"
    "encoding/json"
    "fmt"
    "net/http"
    "github.com/lib/pq"
)

type SensorData struct {
    SensorID     int     `json:"sensor_id"`
    ApproachID   int     `json:"approach_id"`
    VehicleCount int     `json:"vehicle_count"`
    AvgSpeed     float64 `json:"avg_speed"`
    Timestamp   string  `json:"timestamp"`
}

func main() {
    db, _ := sql.Open("postgres", 
        "user=dheerajmurthy dbname=traffic_db password=dheeraj")
    defer db.Close()

    http.HandleFunc("/api/sensors", func(w http.ResponseWriter, r *http.Request) {
        rows, _ := db.Query("SELECT * FROM sensor_data ORDER BY timestamp DESC LIMIT 10")
        // Encode and return JSON...
    })

    fmt.Println("Server starting on :8080")
    http.ListenAndServe(":8080", nil)
}
\end{lstlisting}

\section{Running the System}

\subsection{Setup}

\begin{enumerate}
    \item PostgreSQL running with traffic\_db
    \item C++20 compiler
    \item Go 1.25+
\end{enumerate}

\subsection{Database Setup}

\begin{lstlisting}[language=bash,style=mystyle]
createdb traffic_db
psql -d traffic_db -f sql/db_create.sql
psql -d traffic_db -f sql/db_data.sql
\end{lstlisting}

\subsection{Build and Run Watcher}

\begin{lstlisting}[language=bash,style=mystyle]
cd watcher/build
cmake .. && make -j4
./watcher
\end{lstlisting}

Expected output:
\begin{verbatim}
Connecting to traffic_db...
Connected.
Loaded 9 rules (version 1.0)
  - [rule-001] High Traffic Response (priority 10)
  - [rule-002] Low Traffic Optimization (priority 5)
  ...
Listening for events... (Ctrl+C to exit)
\end{verbatim}

Sample event from running system:
\begin{verbatim}
[sensor_update_channel] {"type": "SENSOR_UPDATE", "sensor_id": 1, 
                        "approach_id": 1, "vehicle_count": 18}
Processing sensor update...
  Data: {
    "sensor_id": 1,
    "approach_id": 1,
    "vehicle_count": 18
  }

[junction_event_channel] {"type": "ACTIVE_PHASE_CHANGE", 
                        "junction_id": 1, "phase_id": 1}
Processing junction event...
\end{verbatim}

\section{System Diagram (Architecture)}

\begin{verbatim}
+----------+     +----------+     +------------+
| Sensor   |---->| PostgreSQL|---->| C++ Watcher|
| Devices  |     |  DB      |     | (Listener) |
+----------+     +----------+     +------------+
                                  |
                                  v (NOTIFY)
                             +-----------+
                             | Dispatcher|
                             +-----------+
                                  |
                    +-------------+-------------+
                    |             |             |
              +-----+-----+ +-----+-----+ +-----+-----+
              |Sensor   | |Signal   | |Junction |
              |Handler  | |Handler | |Handler |
              +---------+ +---------+ +---------+
                    |
                    v
              +-----------+
              | ECA Rules |
              +-----------+
                    |
                    v
              +----------------+
              | Action Execute|
              +----------------+
\end{verbatim}

\section{Future Improvements}

\begin{itemize}
    \item Connect handlers to ECA rules for automatic response
    \item Add time-window aggregation for sensor data
    \item Implement adaptive signal timing algorithms
    \item Add Prometheus metrics endpoint
    \item WebSocket API for real-time dashboard
    \item Unit tests for handlers
\end{itemize}

\section{Conclusion}

This traffic management system demonstrates event-driven architecture using PostgreSQL LISTEN/NOTIFY. The modular design allows independent scaling of components while maintaining real-time responsiveness.

\bigskip
\textbf{Technologies Used:}
\begin{itemize}
    \item PostgreSQL 17 (NOTIFY, triggers, stored functions)
    \item C++20 (watcher with pugixml, nlohmann-json)
    \item Go 1.25 (REST API server)
    \item CMake (build system)
\end{itemize}

\end{document}